
\section{Architecture Integration}

\subsection{Model Support Matrix}

Adaptive SVD supports the following attention variants:

\begin{table}[H]
\centering
\small
\begin{tabular}{lllllll}
\toprule
\textbf{Architecture} & \textbf{Class} & \textbf{Nested?} & \textbf{RoPE?} & \textbf{GQA?} & \textbf{T5 bias?} & \textbf{Verified} \\
\midrule
DistilBERT & \texttt{MultiHeadSelfAttention} \& No \& No \& No \& No &  \\
RoBERTa & \texttt{BertSelfAttention} \& Yes \& No \& No \& No &  \\
GPT-2 & \texttt{GPT2Attention} \& No \& No \& No \& No &  \\
Qwen3.5 & \texttt{Qwen3Attention} \& No \& Yes \& Yes \& No \& partial \\
T5 & \texttt{T5Attention} \& No \& No \& No & \textbf{Yes} \& not tested \\
\bottomrule
\end{tabular}
\end{table}

\subsection{Weight Extraction}

\subsubsection{DistilBERT (Flat)}

Direct extraction:

\begin{lstlisting}
W\_Q = src.q\_lin.weight  # Linear(768, 768)
W\_K = src.k\_lin.weight
W\_V = src.v\_lin.weight
W\_O = src.out\_lin.weight
\end{lstlisting}

All four projections are inside \texttt{MultiHeadSelfAttention}. \texttt{is\_nested\_encoder = False}.

\subsubsection{RoBERTa (Nested)}

Extract from inner \texttt{BertSelfAttention}:

\begin{lstlisting}
W\_Q = src.self.query.weight
W\_K = src.self.key.weight
W\_V = src.self.value.weight
# W\_O is not extracted -- outer BertSelfOutput.dense applies it
\end{lstlisting}

\texttt{is\_nested\_encoder = True}. During forward, the compressed \texttt{context} is 
returned without applying \texttt{out\_proj}; the outer \texttt{BertSelfOutput.dense} 
applies its own (unchanged) projection.

\textbf{Consequence.} In RoBERTa, the output projection $W_O$ is \textbf{not compressed}, 
only $W_Q, W_K, W_V$.

\subsubsection{GPT-2 (Fused Conv1D)}

GPT-2 fuses Q, K, V into a single \texttt{c\_attn} Conv1D:

\begin{lstlisting}
fused = src.c\_attn.weight.T  # (3*768, 768) after transpose
W\_Q, W\_K, W\_V = fused.chunk(3, dim=0)
W\_O = src.c\_proj.weight.T
\end{lstlisting}

The transpose handles the HuggingFace Conv1D layout: it stores weights as 
\texttt{(in\_features, out\_features)} instead of the standard \texttt{(out\_features, in\_features)}.

\subsubsection{T5 (Relative Attention Bias)}

T5 uses \texttt{relative\_attention\_bias} (an \texttt{nn.Embedding}) instead of learned 
position embeddings:

\begin{lstlisting}
self.relative\_attention\_bias = src.relative\_attention\_bias
self.has\_relative\_attention\_bias = True
\end{lstlisting}

The bias is preserved as a separate module and added to attention scores 
in forward. The bucket function handles bidirectional (encoder) and causal 
(decoder) modes.

\subsubsection{Qwen3.5 (GQA + RoPE)}

Qwen3.5 uses GQA with $H_{kv} < H$:

\begin{lstlisting}
W\_Q = src.q\_proj.weight  # (H * d\_h, d) = (2048, 1024)
W\_K = src.k\_proj.weight  # (H\_kv * d\_h, d) = (512, 1024)
W\_V = src.v\_proj.weight  # (512, 1024)
W\_O = src.o\_proj.weight  # (1024, 2048)
\end{lstlisting}

RoPE is applied after projection:

\begin{lstlisting}
Q, K = apply\_rotary\_pos\_emb(Q, K, cos, sin)
\end{lstlisting}

\subsection{The `AdaptiveSVDLinear` Module}

\begin{lstlisting}
class AdaptiveSVDLinear(nn.Module):
    def \_\_init\_\_(self, in\_features, out\_features, svd\_rank, residual\_rank=16, ...):
        super().\_\_init\_\_()
        self.svd\_rank = min(svd\_rank, in\_features, out\_features)
        self.residual\_rank = min(residual\_rank, in\_features, out\_features)
        
        # Frozen SVD base
        self.U = nn.Parameter(
            torch.empty(out\_features, self.svd\_rank),
            requires\_grad=False
        )
        self.V = nn.Parameter(
            torch.empty(self.svd\_rank, in\_features),
            requires\_grad=False
        )
        
        # Trainable residual
        if self.residual\_rank \textgreater{} 0:
            self.lora\_B = nn.Parameter(
                torch.empty(out\_features, self.residual\_rank)
            )
            self.lora\_A = nn.Parameter(
                torch.empty(self.residual\_rank, in\_features)
            )
        
        if bias:
            self.bias = nn.Parameter(torch.empty(out\_features))
        
        self.reset\_parameters()
    
    def forward(self, x):
        Vx = F.linear(x, self.V)
        out = F.linear(Vx, self.U)
        if self.lora\_A is not None:
            Ax = F.linear(x, self.lora\_A)
            out = out + F.linear(Ax, self.lora\_B)
        if self.bias is not None:
            out = out + self.bias
        return out
\end{lstlisting}

\textbf{Parameter count:} $r_P(d_\text{in} + d_\text{out}) + r_L(d_\text{in} + d_\text{out}) + d_\text{out}$.

\subsection{Injection Algorithm}

\textbf{Algorithm 2:} Inject Adaptive SVD into a full model
\begin{lstlisting}
Input:  model M, config c, hyperparameters $\tau$, r\_L
Output: M with attention replaced

1:  for   in transformer\_layers(M) do
2:      src $\leftarrow$ find\_attention( )                     Old attention
3:      if src is None: continue
4:      new $\leftarrow$ AdaptiveSVDAttention(c, $\tau$, r\_L, pretrained\_attn=src)
5:      new.copy\_weights(src)                        SVD + adapter
6:      replace\_attention( , new)
7: end for
8: return M
\end{lstlisting}

\textbf{Layer paths:}

\begin{table}[H]
\centering
\small
\begin{tabular}{lll}
\toprule
\textbf{Model} & \textbf{Layer path} & \textbf{Attention path} \\
\midrule
DistilBERT & \texttt{distilbert.transformer.layer[i]} & \texttt{.attention} \\
RoBERTa & \texttt{roberta.encoder.layer[i]} & \texttt{.attention.self} \\
GPT-2 & \texttt{transformer.h[i]} & \texttt{.attn} \\
T5 & \texttt{encoder.block[i].layer[0]} & \texttt{.SelfAttention} \\
\bottomrule
\end{tabular}
\end{table}

\subsection{Handling Special Cases}

\subsubsection{T5 Relative Attention Bias}

Forward:

\begin{lstlisting}
position\_bias = kwargs.get("position\_bias", None)
if position\_bias is None and self.has\_relative\_attention\_bias:
    position\_bias = self.compute\_bias(N\_q, N\_k, device=hidden\_states.device)
\end{lstlisting}

The \texttt{compute\_bias} function uses \texttt{\_relative\_position\_bucket} to map relative 
positions to buckets, then looks up the bias embedding.

\subsubsection{GQA Expansion}

\begin{lstlisting}
if self.kv\_repeat \textgreater{} 1:
    K = repeat\_kv(K, self.kv\_repeat)
    V = repeat\_kv(V, self.kv\_repeat)
\end{lstlisting}

\subsubsection{RoPE}

\begin{lstlisting}
if position\_embeddings is not None:
    cos, sin = position\_embeddings
    Q, K = apply\_rotary\_pos\_emb(Q, K, cos, sin)
\end{lstlisting}

\subsubsection{Causal Masking}

\begin{lstlisting}
if is\_causal:
    causal\_tril = torch.tril(torch.ones((N\_q, N\_k), device=...))
    scores = scores.masked\_fill(~causal\_tril.unsqueeze(0).unsqueeze(0), -10000.0)
\end{lstlisting}

\subsection{Registry}

\begin{lstlisting}
@ATTENTION\_REGISTRY.register("adaptive\_svd\_attn")
@ATTENTION\_REGISTRY.register("adaptive\_svd")
@ATTENTION\_REGISTRY.register("adaptive\_svd\_attention")
class AdaptiveSVDAttention(BaseAttention):
    ...
\end{lstlisting}

\subsection{Failure Modes and Fallbacks}

\begin{table}[H]
\centering
\small
\begin{tabular}{lll}
\toprule
\textbf{Failure} & \textbf{Detection} & \textbf{Response} \\
\midrule
Missing weights & \texttt{copy\_weights} raises \& Return None \\
Shape mismatch \& Assertion \& Skip layer \\
Near-zero Frobenius \& Check \textless{} 1e-4 \& Skip layer \\
NaN in SVD & \texttt{torch.isnan} \& Fallback \\
T5 bias mismatch \& Embedding shape \& Skip \\
\bottomrule
\end{tabular}
\end{table}

\subsection{Numerical Precision}

\begin{table}[H]
\centering
\small
\begin{tabular}{ll}
\toprule
\textbf{Component} & \textbf{Precision} \\
\midrule
SVD \& FP32 \\
Factor storage \& Model dtype \\
Forward \& Model dtype \\
Adapter init \& FP32 then cast \\
\bottomrule
\end{tabular}
\end{table}

\subsection{Compatibility}

\subsubsection{LoRA Wrappers}

If the source module is a \texttt{LoRALinear} wrapper, Adaptive SVD extracts the 
inner \texttt{base\_layer} for SVD and preserves the outer LoRA parameters:

\begin{lstlisting}
target\_mod = getattr(module, "base\_layer", module)
\end{lstlisting}

\subsubsection{Encoder-Decoder}

For T5, the same \texttt{AdaptiveSVDAttention} is used in both encoder and decoder 
blocks. The decoder side has \texttt{is\_decoder = True}, which enables causal masking.